\appendix

\section{Formalization map and publication gates}
\label{app:formalization-status-v3}

This appendix records the relation between the manuscript and the private Lean
4.31 project. It is part of the verification draft because the size of the
formal code base must not be confused with completion of the top-level theorem.
A declaration is treated as proof evidence only after exact-signature review,
warning-fatal replay under the pinned toolchain, a forbidden-shortcut scan, an
axiom audit, and independent mathematical review.

\subsection{Status vocabulary}

\begin{description}
\item[Welded.] The exact declaration has passed focused and root replay, its
axioms have been audited, and it belongs to the canonical private import
closure.
\item[Running.] An exact one-hole theorem request has been approved and
submitted, but no returned candidate is proof evidence until the complete
independent replay protocol has passed.
\item[Needs review.] The mathematical task is known, but its exact theorem
signature, constants, quantifiers, and dependency ledger have not all been
frozen.
\item[Blocked.] A downstream theorem whose declared dependencies are not yet
all welded.
\end{description}

\subsection{Section 8 finite chain}

The following finite Section~8 declarations are welded privately:

\begin{longtable}{>{\raggedright\arraybackslash}p{0.20\linewidth}
                  >{\raggedright\arraybackslash}p{0.52\linewidth}
                  >{\raggedright\arraybackslash}p{0.18\linewidth}}
\toprule
Manuscript role & Lean declaration or DAG node & Status\\
\midrule
\endhead
High-support encoding &
\texttt{encodeRawNearMiddle\_injective} & Welded\\
Physical fiber and partial weight &
\texttt{profileHighSkeletonWeight\_eq\_fourEndpointPartialAggregateWeight} & Welded\\
One global deficit charge &
\texttt{profileHighSkeletonWeight\_le\_fourEndpointFullReference\_mul\_deficitCharge} & Welded\\
Fixed-support deficit product &
\texttt{sum\_fourEndpointAllHighChoiceWeight\_le\_threeQuarterProduct} & Welded\\
Encoded full-support comparison &
\texttt{profileHighSkeletonWeight\_le\_fourEndpointEncodedFullSupportCharge} & Welded\\
Weighted realized-table regrouping &
\texttt{sum\_fourEndpointFullSupportReferenceWeight\_mul\_tableWeight\_eq\_sum\_realized\_W\_mul\_tableWeight} & Welded\\
Realized-table deficit sum &
\texttt{sum\_profileHighSkeletonWeight\_le\_realizedTableThreeQuarterProduct} & Welded\\
Coarse phase corridor &
\texttt{eventually\_five\_fourths\_log\_sub\_le\_q\_mul\_endpointBudget} & Welded\\
Canonical common-charge smallness &
\texttt{eventually\_fourEndpointThreeQuarterRho\_le\_one} & Running\\
\bottomrule
\end{longtable}

The welded finite chain verifies the completion-free physical-fiber
factorization, the exact local-to-full comparison, the single global
falling-factorial loss, the deficit product, and both unweighted and weighted
regrouping over realized endpoint tables. The current running theorem is the
remaining real-to-\texttt{ENNReal} phase estimate for the explicit sum of the
sixteen endpoint bases.

The main text uses the coarse common-charge route
\[
  \rho_{16}\le1
  \quad\Longrightarrow\quad
  \left(1+(\alpha+1)\rho_{16}\right)^K.
\]
This route is aligned with the private finite reduction. The stronger
pointwise hypothesis $\rho_{ij}\le1/2$ yields the table-preserving product
$\prod_{i,j}(1+2\rho_{ij})^{L_{ij}}$ and remains a reusable corollary, but it is
not required for the shortest proof closure.

\subsection{Other remaining theorem nodes}

The top-level proof also depends on the following nodes.

\begin{longtable}{>{\raggedright\arraybackslash}p{0.22\linewidth}
                  >{\raggedright\arraybackslash}p{0.49\linewidth}
                  >{\raggedright\arraybackslash}p{0.18\linewidth}}
\toprule
Paper section & Exact remaining obligation & Status\\
\midrule
\endhead
Sections 2--4 & Freeze the concrete phase center, slope, and root package, then
export the chromatic lower-tail theorem with the required full-sequence
quantifiers. & Needs review\\
Section 5 & Assemble the welded four-support entropy certificate, the root
separation, tangent rounding, and the positive signed first moment into one
exact theorem consumed by the final adapter. & Needs review\\
Section 7 & Split and prove the empty-corner, central-rate, and full-corner
asymptotic sums uniformly in the phase; only the exact full-corner reindexing is
currently welded. & Needs review\\
Sections 8--9 & Combine phase smallness, endpoint transport, the skeleton
quotient estimate, and the welded residual attachment theorem into the global
normalized-second-moment hypothesis. & Blocked on the preceding nodes\\
Sections 10--11 & The conditional rare-seed amplifier and final event adapter
are welded, but their concrete hypotheses must be instantiated by the preceding
nodes. & Conditional adapter welded\\
\bottomrule
\end{longtable}

The partial-diagonal node is the largest remaining analytic formalization. It
should be divided into four exact declarations rather than submitted as one
broad request:
\[
  \text{empty corner},\quad
  \text{central rate},\quad
  \text{full-corner wrapper},\quad
  \text{partial-diagonal assembly}.
\]
The same discipline should be used for the phase and first-moment nodes.

\subsection{Publication gate}

The compile-time switch \verb|\ErdosProofClosedtrue| may be enabled only after
all of the following hold on one integrated commit:

\begin{enumerate}
\item the canonical common-charge smallness theorem is welded;
\item the skeleton quotient and global attachment sum are welded;
\item the complete partial-diagonal asymptotic theorem is welded;
\item the concrete phase/chromatic-tail package is welded;
\item the signed four-size first-moment assembly is welded;
\item the conditional final adapter is replayed against those exact outputs;
\item the public theorem-facing Lean root is regenerated and warning-fatal;
\item the manuscript constants, quantifiers, bibliography, and theorem names
are synchronized with that root;
\item the full TeX source compiles without unresolved references, citations,
or status placeholders;
\item an independent mathematical review confirms that no private or public
staging lemma is being counted twice.
\end{enumerate}

Until these gates close, the verification banner and conditional theorem
statement in the front matter are mandatory.

\subsection{Recommended theorem-facing Lean organization}

The internal project contains hundreds of fine-grained modules, including both
legacy and simplified proof routes. For publication, the formalization should
export a small curated root:
\[
\begin{array}{l}
\texttt{Erdos625/Paper/Phase.lean},\\
\texttt{Erdos625/Paper/FirstMoment.lean},\\
\texttt{Erdos625/Paper/PartialDiagonals.lean},\\
\texttt{Erdos625/Paper/HighSkeletons.lean},\\
\texttt{Erdos625/Paper/Attachments.lean},\\
\texttt{Erdos625/Paper/Amplification.lean},\\
\texttt{Erdos625/Paper/Main.lean}.
\end{array}
\]
The paper-facing root should import only the declarations corresponding to
named manuscript results. The complete historical development may remain in a
separate exhaustive root, and the superseded cycle/walk route should be marked
as legacy rather than appearing in the main theorem dependency chain.

A machine-readable manifest should map every named paper result to its exact
Lean declaration, source file, commit, status, and axiom audit. The
formalization appendix and reproducibility table can then be generated from
that manifest, eliminating manual status drift.
